\chapter{Introduction}
\label{ch:introduction}

ATTPCROOTv2 (ATTPCROOT in the rest of this manual) is a C++ framework for the simulation and analysis of
data from the Active Target Time Projection Chamber (AT-TPC) and related detectors. It is built on ROOT and
FairRoot: every processing step is a \code{FairTask} that reads branches of a ROOT tree and writes new ones, and
every analysis is a ROOT macro that assembles those tasks into a run.

This manual describes the framework as it stands on the \code{OpenKF-Claude} branch. It is written for physicists who
use the framework to simulate an experiment, reconstruct its data and extract physics, and for those who extend it.
Every chapter explains what a component does, how it is configured, and why the choices matter, with an example
taken from a macro in the repository.

\section{The detector in brief}

In an active target the counting gas is also the reaction target. A beam enters the gas volume, a reaction takes
place somewhere along its path, and the charged products ionise the gas. The ionisation electrons drift in a uniform
electric field towards the pad plane, where a Micromegas amplifies them. The AT-TPC pad plane has about ten thousand
pads, read out by the GET electronics, which samples each pad's signal in time buckets. The pad position gives the
$x$ and $y$ coordinates of a track point, the drift time gives $z$, and the amplitude gives the deposited charge. A
solenoid field (2.85~T for the a1975 experiment used as the running example) bends the tracks into helices, whose
curvature gives the rigidity $B\rho$ of each particle.

A complete measurement therefore needs: the positions of the ionisation along each track (from the traces, by
pulse-shape analysis), the separation of the cloud of points into tracks (pattern recognition), a fit of each track
to obtain its momentum and angles at the reaction vertex (track fitting), and the identity of the particle (from the
energy loss and the rigidity). ATTPCROOT provides all of these, together with a full simulation that produces the
same kind of data from a model of the experiment.

\section{Processing chain}

Simulated and experimental data follow separate paths up to the pad traces and the same path after them.

\begin{center}
\small
\begin{tabular}{p{0.45\textwidth}p{0.45\textwidth}}
\toprule
Simulation & Experiment \\
\midrule
Event generators (beam, reaction) & \\
Geant4 transport $\rightarrow$ \code{AtMCPoint} & \\
Electron drift and diffusion (\code{AtClusterizeTask}) & \\
Pad signals (\code{AtPulseTask}) $\rightarrow$ \code{AtRawEvent} & Unpacking (\code{AtUnpackTask}) $\rightarrow$ \code{AtRawEvent} \\
\midrule
\multicolumn{2}{c}{Pulse-shape analysis (\code{AtPSAtask}) $\rightarrow$ \code{AtEvent} (hits)} \\
\multicolumn{2}{c}{Optional cleaning, space-charge correction} \\
\multicolumn{2}{c}{Pattern recognition (\code{AtPRAtask}) $\rightarrow$ \code{AtPatternEvent} (tracks)} \\
\multicolumn{2}{c}{Particle identification (\code{AtPIDTask}) $\rightarrow$ \code{AtPIDEvent}} \\
\multicolumn{2}{c}{Track fitting (\code{AtFitterTask}) $\rightarrow$ \code{AtTrackingEvent}} \\
\multicolumn{2}{c}{Analysis macros: kinematics, excitation energy, cross sections} \\
\bottomrule
\end{tabular}
\end{center}

Because the reconstruction is shared, a simulated experiment can be processed with exactly the same tasks and
settings as the data. This is how acceptances and resolutions are measured, and it is only meaningful when the
simulation reproduces the detector response (gain, noise, saturation, pad response) closely enough. Chapters
\ref{ch:simulation} and \ref{ch:unpacking} point out where the two paths differ.

\section{Modules}

The source tree is organised in CMake libraries, one per directory. Each produces a shared library \code{libAt...}
with a ROOT dictionary, so its classes can be used from macros.

\begin{center}
\small
\begin{tabular}{p{0.2\textwidth}p{0.72\textwidth}}
\toprule
Module & Role \\
\midrule
\code{AtData} & persisted event, hit, pad, track and fit containers (Chapter~\ref{ch:datamodel}) \\
\code{AtSimulationData} & Monte Carlo truth: points, tracks, stack, vertex propagation \\
\code{AtGenerators} & beam and reaction generators (Chapter~\ref{ch:simulation}) \\
\code{AtDetectors} & detector geometries, sensitive volumes, field (Chapter~\ref{ch:parameters}) \\
\code{AtDigitization} & electron drift, pad signals, trigger (Chapter~\ref{ch:simulation}) \\
\code{AtMap} & pad-plane maps and pad inhibition (Chapter~\ref{ch:parameters}) \\
\code{AtParameter} & run-time parameter containers (Chapter~\ref{ch:parameters}) \\
\code{AtUnpack} & raw data (GRAW, HDF5, FRIBDAQ) to \code{AtRawEvent} (Chapter~\ref{ch:unpacking}) \\
\code{AtReconstruction} & PSA, filters, cleaning, pattern recognition, PID and fitting tasks
  (Chapters~\ref{ch:unpacking}--\ref{ch:pid}) \\
\code{AtTools} & energy loss, propagation, PID estimators, sampling, cleaning algorithms (Chapter~\ref{ch:tools}) \\
\code{AtAnalysis}, \code{AtS800} & higher-level analysis and the S800 spectrometer (Chapter~\ref{ch:tools}) \\
\code{AtEventDisplay} & the event viewer (Chapter~\ref{ch:visualization}) \\
\bottomrule
\end{tabular}
\end{center}

Besides the libraries, the repository holds the run-time assets: \code{geometry/} (detector geometries and media),
\code{parameters/} (parameter files), \code{scripts/} (pad maps and lookup tables), \code{resources/} (energy-loss
tables, correction maps), and \code{macro/} (example, test and experiment macros).

\section{How to read this manual}

Chapter~\ref{ch:installation} explains how to build the framework and run macros. Chapters
\ref{ch:datamodel}--\ref{ch:visualization} describe the framework component by component, in the order in which
data flow through it. Chapter~\ref{ch:workflows} assembles them into complete simulation and data workflows, using
the a1975 experiment ($^{16}$C on a deuterium target) as the worked example. Chapter~\ref{ch:conventions} collects the
coordinate, sign and unit conventions, which are the most common source of silent errors. Chapter~\ref{ch:development}
covers contributing to the framework.

Class, file, task and branch names are set in \code{typewriter}. Setters are listed with the default value taken
from the code. Lengths are in millimetres, energies in MeV, fields in tesla and rigidities in tesla-metres unless
stated otherwise.
